\chapter{Skyhooks and Orbital Transfer}
\label{ch:23}

The fifth part of the book turns from the Earth to the space around it, and it begins with the concept that is the best developed of the orbital proposals and the one on which the others in this part lean. A skyhook, in the sense used here, is a long tether in orbit that rotates about its center of mass so that its lower tip, at the lowest point of each revolution, moves relative to the Earth's surface at a speed much less than the orbital speed, and its upper tip, half a revolution later, moves at a speed much greater. A payload launched to meet the lower tip at the modest speed of an accelerator or a small rocket can be caught there and carried around to the upper tip, from which it is released at a speed that takes it to a higher orbit or on an escape trajectory. The concept was developed by Moravec in 1977~\cite{moravec1977} as a non-synchronous orbital skyhook and was elaborated in the following decades by Forward, Hoyt, and others, and a large literature of engineering analysis exists for it, which is why this program, among those of the book, has the strongest claim to the second readiness class. The author's notes combine the skyhook with the geothermal mass accelerator of Chapter 21 in an assembly in which the accelerator supplies the first part of the velocity and the skyhook the remainder, and the combination is the subject of the present chapter.

The first thing to understand about the skyhook is what it does and does not do. It does not create energy. The payload's gain of energy, and of momentum, is drawn from the orbital motion of the tether and its hub, which slows and drops, and unless the loss is made good the system decays after a few launches. The restoration of the hub's momentum can be accomplished in three ways. Electric propulsion can expend propellant to raise the orbit again; an electrodynamic tether can push against the magnetic field of the Earth, exchanging the hub's momentum with that of the planet's field at the expense of electrical power; and traffic in the opposite direction can return momentum to the hub, as when a payload descending from a higher orbit is caught and slowed. The third is the most elegant, since a system in which the mass launched upward is matched by mass returned downward conserves the hub's momentum on average and requires no external supply. It places a requirement on the economy that it served: the traffic must be balanced, or the imbalance must be paid for. The propellant cost of the first method is calculable and is not trivial. With an electric thruster of exhaust velocity $29.4$ km/s, corresponding to a specific impulse of $3000$ s, restoring the hub after a payload has gained $3.5$ km/s of velocity requires propellant equal to $12$ per cent of the payload's mass, and after a gain of $7$ km/s, $24$ per cent. These fractions are small against the order-one fractions of rocket propellant, which is the reason the concept is attractive, but they are not zero, and the supply of propellant to the hub is a continuing logistic obligation.

The second point concerns the tether. The strength of a material is measured in the length of a column of it that would just support itself, or equivalently in the energy per unit mass at which the material fails, which for steel is about $0.1$ MJ per kilogram, for aramid fibers about $2.5$, for polyethylene and polybenzoxazole fibers about $3.5$ to $3.7$, and for carbon nanotubes, in the best laboratory samples of short length, many times more, though the strength of macroscopic yarns has been far below that of individual tubes. The strength per unit mass determines a characteristic velocity, $V_c=\sqrt{2\sigma/\rho}$, where $\sigma$ is the working stress and $\rho$ the density, and the mass of a rotating tether depends on the ratio of the tip velocity to this characteristic velocity through a function that increases extremely rapidly. At small ratios the tether is light, a modest multiple of the payload; at ratios above one, its mass grows faster than exponentially, and a tether of any practical mass can be built only if the tip velocity is kept comfortably below the characteristic velocity. Taking a working stress equal to half the tensile strength, the characteristic velocity of aramid is $1.6$ km/s, of polybenzoxazole $1.9$ km/s, and of a hypothetical nanotube yarn with a strength of $10$ GPa $2.8$ km/s. A tether with a tip velocity of $2$ km/s then has a mass of ten times the payload for aramid, $4.6$ times for polybenzoxazole, and $1.5$ times for the hypothetical yarn. At a tip velocity of $3.5$ km/s the ratios are $526$, $86$, and $10$. The qualitative lesson is the one that the history of the concept has taught: the skyhook is feasible with materials of present strength at tip velocities of one to two kilometers per second, which is enough to reduce the burden on the lower stages of a launch but not enough to replace them, and it becomes a transportation system for escape trajectories only if materials of far greater strength become available. The assembly of the author's notes, in which an accelerator supplies several kilometers per second at the bottom, is a way of reducing the demand on the tether, and the quantitative case for the combination is that it moves the requirement from the material, which is unavailable, to the accelerator, which is merely very difficult.

The geometry of capture is the third point. The lower tip of the tether passes through the lowest point of its trajectory at a speed relative to the ground equal to the orbital speed minus the tip velocity, and the payload must arrive at the right place at the right time, within a window that is short because the tip moves fast. A tether of $100$ km length rotating with a tip velocity of $3.5$ km/s has a period of about three minutes. The payload must be guided to the tip with precision of meters, over a distance of hundreds of kilometers, by active guidance, and any error in the arrival is either corrected by the payload's own maneuvering or ends in a missed encounter or a collision. The loads on the payload at capture are those of an acceleration of the order of ten times gravity for a tether of this size, which is acceptable for freight and unacceptable for passengers, and a gentler arrangement requires a longer tether at the same tip velocity: the acceleration is $V^2/L$, so doubling the length halves it. A tether of several hundred kilometers is therefore the natural scale for a system that is to carry people, and a tether of that length multiplies the hazards of the environment through which it passes.

The fourth point is survival. A long thin structure in low Earth orbit is exposed to micrometeoroids and to orbital debris, and the probability that it will be struck by an object large enough to cut it, over the lifetime of the system, is not negligible. The surviving designs avoid single-strand failure by multi-line construction, in which the load is carried by a network of lines with redundant cross-links so that a cut in one line does not part the whole, and the network is repaired by robotic agents that traverse it. The survival of the structure over its design life then depends on the rate of cuts, the capacity of the network to tolerate them, and the rate of repair, and the analysis, like that of redundancy in Chapter 22, depends on whether the events are independent. A collision with a large piece of debris may cut many lines at once, and the growth of the debris population in the orbits of interest, which is a consequence of past launches, is a common cause of exactly the kind that defeats redundancy. The orbital regime in which the tether must live is a commons in Ostrom's sense, and its governance is a precondition of the engineering.

The last point concerns orbital transfer itself, which is the subject of the second half of the chapter's title. The transfer between circular orbits of different radii by the most economical two-impulse maneuver, the Hohmann transfer~\cite{hohmann1925}, requires a total velocity increment that for a transfer from low Earth orbit at $400$ km altitude to geostationary orbit is $3.86$ km/s, composed of $2.40$ km/s at the start and $1.46$ km/s at arrival, and takes $5.3$ hours. A skyhook that adds $2V_t$ to the payload's velocity with a tip velocity of $2$ km/s supplies $4$ km/s, which exceeds this requirement, so that a single tether of modest tip velocity, suitably placed, could deliver payloads from low orbit to geostationary transfer without propellant. This is the application for which the concept has been most seriously studied, and it is a more modest and more credible program than the launch of payloads from the ground. It turns the skyhook into infrastructure for traffic among orbits and not a means of reaching orbit, and in that role it connects to the orbital elevator of the next chapter, which would provide the reverse link.

The assessment under the standard of Chapter 16 follows. The mechanism is explicit and the energy balance is written out. The material limits are known, and they confine the near-term concept to tip velocities of one to two kilometers per second. The dependencies are the accelerator or rocket that supplies the payload to the tether, the guidance system, the orbital environment, and the supply of propellant for the hub. The concept is placed in the second readiness class for transfer between orbits, and in the third for the assembly with the accelerator reaching escape velocity. It has no bearing on the satisfaction of present obligations, and its development, if it proceeds, should be funded as a technology demonstration, beginning with a short tether that can validate the dynamics and the capture, and proceeding only upon successful demonstration.

\section*{Formalization}

Consider a uniform-rotation tether of length $L$ (from hub to tip) carrying a tip mass $m_p$, rotating at angular velocity $\omega$, with tip velocity $V=\omega L$ relative to the hub. Let the material have working stress $\sigma$ and density $\rho$, and let the cross-sectional area $A(r)$ be tapered so that the stress equals $\sigma$ everywhere. Define $V_c=\sqrt{2\sigma/\rho}$ and $x=V/V_c$.

\begin{proposition}
The tether mass (one arm) is $M_t=m_p\sqrt\pi\,x\,e^{x^2}\mathrm{erf}(x)$. The mass ratio $M_t/m_p$ is increasing in $x$, behaves as $\tfrac{2}{\sqrt{\pi}}\sqrt\pi x^2=2x^2$ for small $x$, and grows as $\sqrt\pi\,x\,e^{x^2}$ for large $x$.
\end{proposition}

\begin{proof}
Equilibrium of a tether element gives $\sigma\,dA/dr=-\rho\omega^2rA$, hence $A(r)=A_0e^{-kr^2}$ with $k=\rho\omega^2/(2\sigma)$. The tip condition is $\sigma A(L)=m_p\omega^2L$. The mass is $M_t=\rho\int_0^LA(r)\,dr=\rho A_0\frac{\sqrt\pi}{2\sqrt k}\mathrm{erf}(\sqrt kL)$, and $A_0=A(L)e^{kL^2}$. Since $\sqrt kL=x$, $kL^2=x^2$ and $\rho\omega^2L/\sigma=2x^2/L$,
\[
\frac{M_t}{m_p}=\frac{\rho\omega^2L}{\sigma}e^{x^2}\,\frac{\sqrt\pi\,L}{2x}\,\mathrm{erf}(x)=\sqrt\pi\,x\,e^{x^2}\,\mathrm{erf}(x).
\]
For small $x$, $\mathrm{erf}(x)\approx 2x/\sqrt\pi$ gives $2x^2$; for large $x$, $\mathrm{erf}(x)\to1$.
\end{proof}

Numerical values follow, with working stress equal to half the tensile strength: aramid ($3.6$ GPa, $1440\ \mathrm{kg\,m^{-3}}$) $V_c=1.58$ km/s; polybenzoxazole ($5.8$ GPa, $1560$) $V_c=1.93$ km/s; a hypothetical yarn of $10$ GPa and $1300\ \mathrm{kg\,m^{-3}}$, $V_c=2.77$ km/s. For $V=1,2,3,3.5$ km/s the mass ratios $M_t/m_p$ are, for aramid, $1.05$, $10.3$, $122$, $526$; for polybenzoxazole, $0.65$, $4.6$, $30$, $86$; for the hypothetical yarn, $0.28$, $1.5$, $5.4$, $10$. The centripetal acceleration at the tip is $V^2/L$, equal to $122\ \mathrm{m\,s^{-2}}=12.5\,g$ for $V=3.5$ km/s and $L=100$ km, with a tip tension for a $1000$ kg payload of $122.5$ kN, and a tip cross-section of $61\ \mathrm{mm^2}$ (a diameter of $8.8$ mm if circular) at a working stress of $2$ GPa.

For the hub's momentum, let a hub of mass $M_h\gg m_p$ impart a velocity gain $\Delta v_p$ to a payload of mass $m_p$.

\begin{proposition}
By conservation of momentum, the hub's velocity decrement is $\Delta v_h=m_p\Delta v_p/M_h$. Restoring it with a thruster of exhaust velocity $v_e$ requires propellant mass $m_{\mathrm{prop}}=M_h(1-e^{-\Delta v_h/v_e})\approx m_p\Delta v_p/v_e$, so the propellant fraction is $m_{\mathrm{prop}}/m_p\approx\Delta v_p/v_e$, independent of $M_h$. If an equal mass descends and is slowed by the same amount $\Delta v_p$, the net decrement is zero.
\end{proposition}

\begin{proof}
Momentum conservation gives $M_h\Delta v_h=m_p\Delta v_p$. The rocket equation~\cite{tsiolkovsky1903} gives the propellant needed to change the hub's velocity by $\Delta v_h$, with $\Delta v_h/v_e\ll1$ allowing the linearization. For the symmetric case, the momenta are equal and opposite.
\end{proof}

With $v_e=29.4$ km/s, $\Delta v_p=2$, $3.5$, $7$ km/s gives $m_{\mathrm{prop}}/m_p=6.8$, $11.9$, $23.8$ per cent.

For orbital transfer, with $\mu=3.986\times10^5\ \mathrm{km^3\,s^{-2}}$, $r_1=6771$ km and $r_2=42164$ km, the transfer ellipse has semimajor axis $a=(r_1+r_2)/2=24467$ km; the circular speeds are $7.673$ and $3.075$ km/s; the transfer speeds at perigee and apogee are $\sqrt{\mu(2/r_1-1/a)}=10.072$ km/s and $\sqrt{\mu(2/r_2-1/a)}=1.618$ km/s, so $\Delta v_1=2.399$ km/s, $\Delta v_2=1.457$ km/s, total $3.857$ km/s, and the transfer time is $\pi\sqrt{a^3/\mu}=5.29$ h.
